

\section{Full optimizer results}
\label{app:master}

Table~\ref{tab:master} expands the compact gain comparison in
Table~\ref{tab:gain-by-task}. It reports the observed range across replicate
runs and two properties of the resulting harnesses for every core contestant;
the two additional harnesses evaluated on GAIA are included here as well. The
generational-ladder runs remain in Figure~\ref{fig:ladder}, where release order
is the comparison of interest.

% Auto-generated by he-bench-paper/tables/T02_master/render.py -- do not edit by hand.
\begin{table*}[t]
\renewcommand{\arraystretch}{0.7}
\centering\small
\setlength{\tabcolsep}{4pt}
\begin{tabular}{@{}ll r r r@{}}
\toprule
\textbf{Model} & \textbf{Harness} & \textbf{Gain} & \textbf{Levers} & \textbf{Tgt tokens (M)}\\
\midrule
\multicolumn{5}{@{}l}{\textbf{OfficeQA}\quad{\footnotesize resolution band $\pm0.045$ (normalized gain), baseline $0.341$}}\\
\midrule
claude-opus-5 & \texttt{claude-code} & $0.59$ {\tiny $(0.54\text{--}0.64)$} & $1.00$ {\tiny $(1.00\text{--}1.00)$} & $9.97$ {\tiny $(3.05\text{--}16.89)$}\\
claude-opus-5 & \texttt{opencode} & $0.63$ {\tiny $(0.60\text{--}0.67)$} & $1.00$ {\tiny $(1.00\text{--}1.00)$} & $3.77$ {\tiny $(3.36\text{--}4.18)$}\\
claude-sonnet-5 & \texttt{claude-code} & $0.53$ {\tiny $(0.51\text{--}0.55)$} & $0.75$ {\tiny $(0.62\text{--}0.88)$} & $5.61$ {\tiny $(4.92\text{--}6.29)$}\\
claude-sonnet-5 & \texttt{opencode} & $0.51$ {\tiny $(0.46\text{--}0.56)$} & $0.94$ {\tiny $(0.88\text{--}1.00)$} & $2.55$ {\tiny $(2.31\text{--}2.79)$}\\
gpt-5.6-sol & \texttt{codex} & $0.49$ {\tiny $(0.47\text{--}0.50)$} & $0.75$ {\tiny $(0.75\text{--}0.75)$} & $2.00$ {\tiny $(1.59\text{--}2.41)$}\\
gpt-5.6-sol & \texttt{opencode} & $0.29$ {\tiny $(0.27\text{--}0.32)$} & $0.56$ {\tiny $(0.50\text{--}0.62)$} & $1.17$ {\tiny $(1.03\text{--}1.32)$}\\
gpt-5.6-terra & \texttt{codex} & $0.07$ {\tiny $(-0.04\text{--}0.17)$} & $0.69$ {\tiny $(0.62\text{--}0.75)$} & $1.45$ {\tiny $(1.31\text{--}1.58)$}\\
gpt-5.6-terra & \texttt{opencode} & $0.14$ {\tiny $(0.00\text{--}0.27)$} & $0.44$ {\tiny $(0.38\text{--}0.50)$} & $1.26$ {\tiny $(1.15\text{--}1.37)$}\\
kimi-k3 & \texttt{kimi-cli} & $0.59$ {\tiny $(0.59\text{--}0.60)$} & $1.00$ {\tiny $(1.00\text{--}1.00)$} & $2.77$ {\tiny $(2.35\text{--}3.18)$}\\
kimi-k3 & \texttt{opencode} & $0.41$ {\tiny $(0.23\text{--}0.59)$} & $0.94$ {\tiny $(0.88\text{--}1.00)$} & $3.00$ {\tiny $(2.21\text{--}3.79)$}\\
\addlinespace[7pt]
\multicolumn{5}{@{}l}{\textbf{BrowseComp-Plus}\quad{\footnotesize resolution band $\pm0.066$ (normalized gain), baseline $0.462$}}\\
\midrule
claude-opus-5 & \texttt{claude-code} & $0.41$ {\tiny $(0.38\text{--}0.44)$} & $1.00$ {\tiny $(1.00\text{--}1.00)$} & $11.78$ {\tiny $(9.95\text{--}13.61)$}\\
claude-opus-5 & \texttt{opencode} & $0.48$ {\tiny $(0.42\text{--}0.55)$} & $1.00$ {\tiny $(1.00\text{--}1.00)$} & $12.50$ {\tiny $(11.13\text{--}13.88)$}\\
claude-sonnet-5 & \texttt{claude-code} & $0.07$ {\tiny $(0.02\text{--}0.12)$} & $0.81$ {\tiny $(0.75\text{--}0.88)$} & $14.47$ {\tiny $(10.58\text{--}18.36)$}\\
claude-sonnet-5 & \texttt{opencode} & $0.15$ {\tiny $(0.11\text{--}0.19)$} & $0.81$ {\tiny $(0.75\text{--}0.88)$} & $16.10$ {\tiny $(13.44\text{--}18.77)$}\\
gpt-5.6-sol & \texttt{codex} & $0.03$ {\tiny $(0.02\text{--}0.03)$} & $0.69$ {\tiny $(0.62\text{--}0.75)$} & $4.04$ {\tiny $(3.33\text{--}4.76)$}\\
gpt-5.6-sol & \texttt{opencode} & $0.09$ {\tiny $(0.05\text{--}0.14)$} & $0.50$ {\tiny $(0.25\text{--}0.75)$} & $7.59$ {\tiny $(5.18\text{--}10.00)$}\\
gpt-5.6-terra & \texttt{codex} & $-0.03$ {\tiny $(-0.06\text{--}0.01)$} & $0.50$ {\tiny $(0.50\text{--}0.50)$} & $7.77$ {\tiny $(7.59\text{--}7.96)$}\\
gpt-5.6-terra & \texttt{opencode} & $0.02$ {\tiny $(0.01\text{--}0.04)$} & $0.56$ {\tiny $(0.50\text{--}0.62)$} & $8.47$ {\tiny $(8.45\text{--}8.50)$}\\
kimi-k3 & \texttt{kimi-cli} & $0.23$ {\tiny $(0.23\text{--}0.23)$} & $1.00$ {\tiny $(1.00\text{--}1.00)$} & $15.65$ {\tiny $(12.38\text{--}18.93)$}\\
kimi-k3 & \texttt{opencode} & $0.16$ {\tiny $(0.09\text{--}0.22)$} & $0.94$ {\tiny $(0.88\text{--}1.00)$} & $13.60$ {\tiny $(10.71\text{--}16.49)$}\\
\addlinespace[7pt]
\multicolumn{5}{@{}l}{\textbf{Terminal-Bench}\quad{\footnotesize resolution band $\pm0.054$ (normalized gain), baseline $0.241$}}\\
\midrule
claude-opus-5 & \texttt{claude-code} & $0.18$ {\tiny $(0.17\text{--}0.20)$} & $0.88$ {\tiny $(0.88\text{--}0.88)$} & $4.03$ {\tiny $(2.71\text{--}5.36)$}\\
claude-opus-5 & \texttt{opencode} & $0.29$ {\tiny $(0.29\text{--}0.29)$} & $0.81$ {\tiny $(0.75\text{--}0.88)$} & $4.46$ {\tiny $(3.86\text{--}5.07)$}\\
claude-sonnet-5 & \texttt{claude-code} & $0.10$ {\tiny $(0.02\text{--}0.18)$} & $0.62$ {\tiny $(0.50\text{--}0.75)$} & $3.19$ {\tiny $(1.75\text{--}4.62)$}\\
claude-sonnet-5 & \texttt{opencode} & $0.15$ {\tiny $(0.09\text{--}0.22)$} & $0.62$ {\tiny $(0.50\text{--}0.75)$} & $4.81$ {\tiny $(2.66\text{--}6.95)$}\\
gpt-5.6-sol & \texttt{codex} & $0.12$ {\tiny $(0.11\text{--}0.12)$} & $0.50$ {\tiny $(0.25\text{--}0.75)$} & $1.79$ {\tiny $(1.44\text{--}2.14)$}\\
gpt-5.6-sol & \texttt{opencode} & $0.13$ {\tiny $(0.12\text{--}0.15)$} & $0.38$ {\tiny $(0.25\text{--}0.50)$} & $1.46$ {\tiny $(1.24\text{--}1.67)$}\\
gpt-5.6-terra & \texttt{codex} & $0.01$ {\tiny $(-0.02\text{--}0.05)$} & $0.50$ {\tiny $(0.38\text{--}0.62)$} & $1.12$ {\tiny $(1.09\text{--}1.14)$}\\
gpt-5.6-terra & \texttt{opencode} & $0.04$ {\tiny $(0.00\text{--}0.09)$} & $0.44$ {\tiny $(0.38\text{--}0.50)$} & $1.72$ {\tiny $(1.13\text{--}2.32)$}\\
kimi-k3 & \texttt{kimi-cli} & $0.16$ {\tiny $(0.13\text{--}0.20)$} & $0.75$ {\tiny $(0.75\text{--}0.75)$} & $3.00$ {\tiny $(2.47\text{--}3.54)$}\\
kimi-k3 & \texttt{opencode} & $0.12$ {\tiny $(0.04\text{--}0.21)$} & $0.81$ {\tiny $(0.75\text{--}0.88)$} & $2.59$ {\tiny $(0.87\text{--}4.32)$}\\
\addlinespace[7pt]
\multicolumn{5}{@{}l}{\textbf{GAIA}\quad{\footnotesize resolution band $\pm0.035$ (normalized gain), baseline $0.000$}}\\
\midrule
claude-opus-5 & \texttt{claude-code} & $0.42$ {\tiny $(0.34\text{--}0.50)$} & $1.00$ {\tiny $(1.00\text{--}1.00)$} & $1.79$ {\tiny $(1.44\text{--}2.13)$}\\
claude-opus-5 & \texttt{goose} & $0.31$ {\tiny $(0.30\text{--}0.32)$} & $0.88$ {\tiny $(0.88\text{--}0.88)$} & $0.42$ {\tiny $(0.30\text{--}0.54)$}\\
claude-opus-5 & \texttt{mini-swe-agent} & $0.36$ {\tiny $(0.34\text{--}0.38)$} & $0.00$ {\tiny $(0.00\text{--}0.00)$} & $1.42$ {\tiny $(0.92\text{--}1.91)$}\\
claude-opus-5 & \texttt{opencode} & $0.47$ {\tiny $(0.40\text{--}0.54)$} & $0.94$ {\tiny $(0.88\text{--}1.00)$} & $1.68$ {\tiny $(0.55\text{--}2.80)$}\\
claude-sonnet-5 & \texttt{claude-code} & $0.33$ {\tiny $(0.23\text{--}0.42)$} & $0.75$ {\tiny $(0.75\text{--}0.75)$} & $0.58$ {\tiny $(0.29\text{--}0.88)$}\\
claude-sonnet-5 & \texttt{goose} & $0.27$ {\tiny $(0.23\text{--}0.31)$} & $0.69$ {\tiny $(0.62\text{--}0.75)$} & $0.29$ {\tiny $(0.23\text{--}0.36)$}\\
claude-sonnet-5 & \texttt{mini-swe-agent} & $0.20$ {\tiny $(0.12\text{--}0.27)$} & $0.00$ {\tiny $(0.00\text{--}0.00)$} & $0.19$ {\tiny $(0.04\text{--}0.33)$}\\
claude-sonnet-5 & \texttt{opencode} & $0.25$ {\tiny $(0.24\text{--}0.26)$} & $0.81$ {\tiny $(0.75\text{--}0.88)$} & $0.35$ {\tiny $(0.22\text{--}0.47)$}\\
gpt-5.6-sol & \texttt{codex} & $0.49$ {\tiny $(0.47\text{--}0.52)$} & $0.69$ {\tiny $(0.62\text{--}0.75)$} & $0.55$ {\tiny $(0.20\text{--}0.91)$}\\
gpt-5.6-sol & \texttt{goose} & $0.20$ {\tiny $(0.13\text{--}0.27)$} & $0.69$ {\tiny $(0.62\text{--}0.75)$} & $0.24$ {\tiny $(0.21\text{--}0.27)$}\\
gpt-5.6-sol & \texttt{mini-swe-agent} & $0.30$ {\tiny $(0.23\text{--}0.36)$} & $0.00$ {\tiny $(0.00\text{--}0.00)$} & $0.24$ {\tiny $(0.13\text{--}0.35)$}\\
gpt-5.6-sol & \texttt{opencode} & $0.31$ {\tiny $(0.28\text{--}0.34)$} & $0.81$ {\tiny $(0.75\text{--}0.88)$} & $0.39$ {\tiny $(0.19\text{--}0.58)$}\\
gpt-5.6-terra & \texttt{codex} & $0.30$ {\tiny $(0.26\text{--}0.35)$} & $0.69$ {\tiny $(0.62\text{--}0.75)$} & $0.42$ {\tiny $(0.17\text{--}0.67)$}\\
gpt-5.6-terra & \texttt{goose} & $0.17$ {\tiny $(0.12\text{--}0.22)$} & $0.81$ {\tiny $(0.75\text{--}0.88)$} & $0.21$ {\tiny $(0.15\text{--}0.27)$}\\
gpt-5.6-terra & \texttt{mini-swe-agent} & $0.16$ {\tiny $(0.14\text{--}0.17)$} & $0.00$ {\tiny $(0.00\text{--}0.00)$} & $0.11$ {\tiny $(0.08\text{--}0.14)$}\\
gpt-5.6-terra & \texttt{opencode} & $0.17$ {\tiny $(0.16\text{--}0.18)$} & $0.62$ {\tiny $(0.62\text{--}0.62)$} & $0.15$ {\tiny $(0.15\text{--}0.15)$}\\
kimi-k3 & \texttt{goose} & $0.24$ {\tiny $(0.00\text{--}0.47)$} & $0.00$ {\tiny $(0.00\text{--}0.00)$} & $0.47$ {\tiny (1 run)}\\
kimi-k3 & \texttt{kimi-cli} & $0.31$ {\tiny $(0.27\text{--}0.34)$} & $0.81$ {\tiny $(0.75\text{--}0.88)$} & $0.37$ {\tiny $(0.35\text{--}0.39)$}\\
kimi-k3 & \texttt{mini-swe-agent} & $0.34$ {\tiny (1 run)} & $0.00$ {\tiny $(0.00\text{--}0.00)$} & $0.86$ {\tiny (1 run)}\\
kimi-k3 & \texttt{opencode} & $0.28$ {\tiny $(0.27\text{--}0.29)$} & $0.75$ {\tiny $(0.62\text{--}0.88)$} & $0.61$ {\tiny $(0.59\text{--}0.62)$}\\
\bottomrule
\end{tabular}
\caption{Optimizer performance by task, as \emph{mean (observed range)} over a contestant's rounds. The Gain column repeats Table~\ref{tab:gain-by-task}; generational-ladder rungs are excluded here and appear in Figure~\ref{fig:ladder}. Rows are ordered by model name, so a contestant sits at the same point in every block; blocks differ in length because two coding harnesses ran on one task only. A parenthetical is the range observed across that contestant's rounds, not a confidence interval --- two rounds do not estimate dispersion. ``--'' is not a zero: it is a measure that was not computable for that cell. In 7 entries the contestant shipped the unmodified seed on at least one round, so that round re-measures the seed rather than a failed edit.}
\label{tab:master}
\end{table*}


% Keep each appendix artifact behind the prose that introduces it and ahead of
% the next appendix section.
\FloatBarrier

\section{The suite}
\label{app:suite}

Table~\ref{tab:suite} gives each task's design constants --- its pinned target
model, its exact split, and its seed agent's held-out baseline --- together with
what the same task scores under each off-the-shelf coding harness on the same
held-out partition. Gain is measured against the seed baseline alone
(Eq.~\ref{eq:normalized-gain}); the off-the-shelf columns are context, not the
reference --- they say what the task's headroom is worth to an existing harness
that nobody optimized.

Scores are means over $K{=}3$ rounds, with $\pm$ denoting the standard error
across round means; timeouts count as zero. GAIA starts from a
non-functional stub. Task data come from GAIA \citep{mialon2023gaia}, OfficeQA
Pro \citep{opsahlong2026officeqaproenterprisebenchmark}, BrowseComp-Plus
\citep{browsecompplus2025}, and Terminal-Bench 2.0
\citep{merrill2026terminalbenchbenchmarkingagentshard}.

% Auto-generated by he-bench-paper/tables/T01_suite/render.py -- do not edit by hand.
\begin{table*}[t]
\centering\footnotesize
\setlength{\tabcolsep}{4pt}
\resizebox{\textwidth}{!}{%
\begin{tabular}{@{}lcrrrrrr@{}}
\toprule
 & & & \multicolumn{5}{c}{\textbf{Off-the-shelf harness}}\\
\cmidrule(l){4-8}
\textbf{Task} & \textbf{Split (d/v/t)} & \textbf{Seed} & \textbf{\texttt{opencode}} & \textbf{\texttt{goose}} & \textbf{\shortstack{\texttt{openhands-}\\\texttt{sdk}}} & \textbf{\shortstack{\texttt{mini-swe-}\\\texttt{agent}}} & \textbf{\shortstack{\texttt{terminus-}\\\texttt{2}}}\\
\midrule
\shortstack[l]{OfficeQA\\[1pt]{\tiny \texttt{deepseek-v4-flash}}} & $49/98/99$ & \shortstack[c]{$0.341$\\[1pt]{\tiny $\pm0.023$}} & \shortstack[c]{$0.727$\\[1pt]{\tiny $\pm0.010$}} & \shortstack[c]{$0.505$\\[1pt]{\tiny $\pm0.129$}} & \shortstack[c]{$0.713$\\[1pt]{\tiny $\pm0.011$}} & \shortstack[c]{$\mathbf{0.734}$\\[1pt]{\tiny $\pm0.009$}} & \shortstack[c]{$0.687$\\[1pt]{\tiny $\pm0.031$}}\\
\shortstack[l]{BrowseComp-Plus\\[1pt]{\tiny \texttt{deepseek-v4-flash}}} & $33/66/66$ & \shortstack[c]{$0.462$\\[1pt]{\tiny $\pm0.020$}} & \shortstack[c]{$0.434$\\[1pt]{\tiny $\pm0.018$}} & \shortstack[c]{$0.615$\\[1pt]{\tiny $\pm0.013$}} & \shortstack[c]{$0.657$\\[1pt]{\tiny $\pm0.035$}} & \shortstack[c]{$\mathbf{0.701}$\\[1pt]{\tiny $\pm0.009$}} & \shortstack[c]{$0.439$\\[1pt]{\tiny $\pm0.023$}}\\
\shortstack[l]{Terminal-Bench\\[1pt]{\tiny \texttt{grok-build}}} & $17/36/36$ & \shortstack[c]{$0.241$\\[1pt]{\tiny $\pm0.009$}} & \shortstack[c]{$\mathbf{0.607}$\\[1pt]{\tiny $\pm0.030$}} & \shortstack[c]{$0.434$\\[1pt]{\tiny $\pm0.024$}} & \shortstack[c]{$0.393$\\[1pt]{\tiny $\pm0.013$}} & \shortstack[c]{$0.364$\\[1pt]{\tiny $\pm0.032$}} & \shortstack[c]{$0.333$\\[1pt]{\tiny $\pm0.016$}}\\
\shortstack[l]{GAIA$^{\S}$\\[1pt]{\tiny \texttt{gpt-5.4-mini}}} & $33/66/66$ & $0.000$ & \shortstack[c]{$0.469$\\[1pt]{\tiny $\pm0.023$}} & \shortstack[c]{$0.207$\\[1pt]{\tiny $\pm0.013$}} & \shortstack[c]{$\mathbf{0.508}$\\[1pt]{\tiny $\pm0.029$}} & \shortstack[c]{$0.172$\\[1pt]{\tiny $\pm0.022$}} & \shortstack[c]{$0.202$\\[1pt]{\tiny $\pm0.005$}}\\
\bottomrule
\end{tabular}}
\caption{The \heb\ suite: each task's pinned target model (second line), its development/validation/test split, its seed agent's held-out score pooled over $K{=}3$ rounds, and what the same task scores under each off-the-shelf coding harness. An off-the-shelf column swaps a stock harness in for the seed program and changes nothing else --- same dataset, same partition, same rounds, same target model --- so the harness is the only variable; best per task in bold. They provide an indication of the seed harness' headroom. Every $\pm$ is the standard error of that cell's three round means ($\mathrm{sd}/\sqrt{K}$): re-run variation, not case-to-case spread. $^{\S}$this seed is a non-functional stub, so its baseline is a measured zero and gain there is the raw held-out score --- the one task on which an optimizer's result and a stock harness's score are directly comparable. Timeouts are scored as zeros rather than dropped per each benchmarks specifications. Task data are drawn from GAIA \citep{mialon2023gaia}, OfficeQA Pro \citep{opsahlong2026officeqaproenterprisebenchmark}, BrowseComp-Plus \citep{browsecompplus2025}, and Terminal-Bench 2.0 \citep{merrill2026terminalbenchbenchmarkingagentshard}.}
\label{tab:suite}
\end{table*}


% Keep the suite table behind its introduction and ahead of the next appendix
% section.  In particular, do not let the later single-column model-effects
% table overtake this double-column table in LaTeX's separate float queues.
\FloatBarrier

\section{Model effects}
\label{app:model-effects}

Table~\ref{tab:model-effects} gives the model effect in gain units, the numeric
form of the ordering Figure~\ref{fig:showdown} plots. Two further estimators of the
same quantity are computed and agree with it exactly on the ordering, so the ranking
does not rest on the additive assumption; only the gain-unit estimator is reported,
because it is the one the body quotes and the only one in units a reader can read
directly as a fraction of headroom.

% Auto-generated by he-bench-paper/tables/T06_model_effects/render.py -- do not edit by hand.
\begin{table}[t]
\centering\small
\setlength{\tabcolsep}{4pt}
\begin{tabular}{@{}l rc@{}}
\toprule
& \multicolumn{2}{c}{\textbf{LSS-$\lambda$}}\\
\cmidrule(l){2-3}
\textbf{Model} & gain units & tier\\
\emph{resolution} & $\pm0.058$ & \\
\midrule
claude-opus-5 & $+0.228$ & 1\\
claude-sonnet-5 & $+0.029$ & 2\\
kimi-k3 & $-0.014$ & 2\\
gpt-5.6-sol & $-0.069$ & 2\\
gpt-5.6-terra & $-0.174$ & 3\\
\bottomrule
\end{tabular}
\caption{\textbf{The model effect on the balanced scope.} LSS-$\lambda$ is the model term of an additive fit over task and model, in normalized-gain units: what swapping the optimizer model is worth once the task is accounted for. This is the numeric form of the right panel of Figure~\ref{fig:showdown}, from the same call, so the two cannot disagree. \emph{resolution} is the estimator's own measured split-round swing and the tier column cuts at it, so a shared tier means ``closer together than re-running the grid moves them'' and ranks within a tier are not resolved. Two further estimators of the same quantity --- gain standardized within task, and mean within-group rank --- are computed but not shown; all three produce the identical ordering (pairwise rank concordance $\geq1.00$), so the ranking here does not depend on the additive assumption. Scope: 15 contestant rows over 3 tasks on \oc\ alone, balanced. GAIA is excluded because its measured-zero baseline makes its gain a raw held-out score rather than a fraction of headroom. No harness effect is reported: this scope holds one harness, so the fit has no harness factor to estimate --- the harness magnitude in Section~\ref{sec:rq1} comes from the paired contrast instead.}
\label{tab:model-effects}
\end{table}


\FloatBarrier

\section{Supporting figures}
\label{app:figures}

All figures use the same run store and analysis set as the body. Where applicable,
points average replicate rounds for one optimizer configuration on one task, and
$\rho$ is Spearman correlation within a task. As in the body, GAIA gain is its raw
held-out score and is not pooled with the other tasks.

\begin{figure*}[t]
\centering
\includegraphics[width=\linewidth]{figures/A10_edit_prevalence.pdf}
\caption{\textbf{Edit prevalence by task.} Share of balanced-grid runs that made
at least one edit in each category; annotations give run counts.
$^{\ddagger}$GAIA starts from a non-functional stub.}
\label{fig:edit-prevalence}
\end{figure*}

% The next artifact is single-column; force the double-column prevalence figure
% to appear first rather than allowing the two float queues to reorder them.
% \FloatBarrier

\begin{figure}[t]
\centering
\includegraphics[width=0.65\linewidth]{figures/A12_role_breadth_vs_gain.pdf}
\caption{\textbf{Shipped edit breadth and gain.} Normalized gain against the
fraction of edit categories present in the submitted harness.
Figure~\ref{fig:levers} reports breadth explored during search.}
\label{fig:role-breadth}
\end{figure}

% \FloatBarrier

% Apaar's R01-R11, rendered through this module against our config and house style
% (see he-bench-paper/figures/apaar_series.py). Placed here wholesale so the pairs can
% be compared on the page; four of them carry the labels of the A figures they replace,
% so no prose reference changed.


\begin{figure*}[t]
\centering
\includegraphics[width=\linewidth]{figures/R04_action_profile.pdf}
\caption{\textbf{Optimizer action profiles.} Share of classified, non-polling
actions by task and optimizer model; polling is shown separately.}
\label{fig:action-profile}
\end{figure*}

\begin{figure*}[t]
\centering
\includegraphics[width=\linewidth]{figures/R05_optimizer_cost_vs_gain.pdf}
\caption{\textbf{Search cost and gain.} Optimizer token use versus normalized
gain; filled points form the non-dominated frontier within each task.}
\label{fig:opt-cost}
\end{figure*}

\begin{figure*}[t]
\centering
\includegraphics[width=\linewidth]{figures/R06_target_cost_vs_gain.pdf}
\caption{\textbf{Shipped-agent cost and gain.} Inference tokens per held-out
case versus normalized gain; filled points form the non-dominated frontier
within each task.}
\label{fig:target-cost}
\end{figure*}

\begin{figure*}[t]
\centering
\includegraphics[width=\linewidth]{figures/R07_edit_share_vs_gain.pdf}
\caption{\textbf{Modification volume and gain.} Lines changed in the target
agent versus normalized gain; task-level rank correlations appear above each
panel.}
\label{fig:edit-share-vs-gain}
\end{figure*}

\begin{figure*}[t]
\centering
\includegraphics[width=0.8\linewidth]{figures/R08_native_vs_common.pdf}
\caption{\textbf{Native versus common optimizer harnesses.} Bars count
model--task pairs favoring the native harness or \oc; color marks a significant
two-sided sign test.}
\label{fig:native-vs-common}
\end{figure*}

\begin{figure*}[t]
\centering
\includegraphics[width=0.7\linewidth]{figures/R09_overfitting.pdf}
\caption{\textbf{Validation versus held-out performance.} Best observed
validation score versus the submitted candidate's test score; the diagonal is
equality. Cells without a validation measurement are omitted.}
\label{fig:overfit}
\end{figure*}

\begin{figure*}[t]
\centering
\includegraphics[width=0.75\linewidth]{figures/R10_harness_sweep.pdf}
\caption{\textbf{GAIA harness sweep.} Mean normalized gain for each evaluated
optimizer-model--harness pairing; dashes denote combinations not run.}
\label{fig:harness-sweep}
\end{figure*}

\begin{figure*}[t]
\centering
\includegraphics[width=\linewidth]{figures/R11_budget_use.pdf}
\caption{\textbf{Case-pass budgets, not evaluation calls, bind.} Median share
of each cap consumed by task and harness scope; annotations count cells that
exhausted a case-pass budget.}
\label{fig:budget-use}
\end{figure*}

% Supporting artifacts belong to the section above.  Flush them before the
% prose-only reproducibility section begins.
\FloatBarrier

\section{Reproducibility}
\label{app:repro}

Each task pins its dataset by immutable reference and its split by a committed
manifest. Splits are exact $20/40/40$ with no overlap, and the split generator
verifies the committed tree byte-for-byte. Baselines are the seed agent's
held-out score pooled over $K{=}3$ rounds and are re-pinned whenever a seed
commit moves, using a script that reuses the original evaluation path. Every
candidate an optimizer produces is an immutable Git commit, dependencies are
pinned by lockfile, and each evaluation runs in an isolated sandbox with
per-case timeouts taken from the dataset's own declared agent clock. The trusted
gateway meters each evaluation's token usage as an input, cached, output, and
total split, and stamps request-log records so that per-trial token attribution
is reproducible from the run artifacts. Every quantity we report about an
optimizer's process is recomputed from its own execution trace and from the
evaluator's record of which evaluations it ran, so the analysis is reproducible
from the released artifacts without re-running any search. The seed agents, split
manifests, build configurations, and
evaluation harness are released with the benchmark.
